\textbf{Задание №1-7} Теоретические сведения.\label{sec:1}

\begin{enumerate}

\item \textbf{Виртуальная память} — механизм, предоставляющий каждому процессу
иллюзию непрерывного изолированного адресного пространства, отображая
виртуальные адреса в физическую память и подкачку на диск через таблицы
страниц.

\textbf{Назначение:}
\begin{enumerate}
    \item абстракция над физической памятью;
    \item изоляция процессов друг от друга;
    \item защита данных;
    \item возможность запускать программы больше объёма RAM;
    \item упрощение управления памятью для приложений.
\end{enumerate}

\textbf{Отличие от физической памяти:}
физическая память — реальные микросхемы RAM с фиксированными адресами;
виртуальная — иллюзия, где каждый процесс видит своё пространство, а MMU
преобразует виртуальные адреса в физические.


\item Механизмы отображения адресов \label{sec:2}

\textbf{MMU} (Memory Management Unit) — аппаратный блок процессора,
преобразующий виртуальные адреса в физические через таблицы страниц и
проверяющий права доступа.

\textbf{Таблицы страниц x86-64 (4 уровня):}
\begin{enumerate}
    \item \textbf{PML4} — Page Map Level 4 (биты 39--47);
    \item \textbf{PDPT} — Page Directory Pointer Table (биты 30--38);
    \item \textbf{PD} — Page Directory (биты 21--29);
    \item \textbf{PT} — Page Table (биты 12--20);
    \item \textbf{offset} — смещение внутри страницы (биты 0--11).
\end{enumerate}

\textbf{LA57} — пятиуровневая адресация с дополнительным PML5
(биты 48--56), расширяет пространство до 128~ПБ.

\textbf{Страница} — минимальная единица (обычно 4~КБ). Размер:
\texttt{getconf PAGE\_SIZE}.

\textbf{TLB} (Translation Lookaside Buffer) — кэш последних трансляций
<<виртуальный~\(\to\)~физический>>. TLB miss вызывает page table walk,
но \emph{не} является page fault.


\item Адресное пространство процесса \label{sec:3}

\begin{enumerate}
    \item \texttt{.text} — машинные инструкции, права \texttt{r-x};
    \item \texttt{.data} — инициализированные глобальные/статические;
    \item \texttt{.bss} — неинициализированные глобальные/статические (нули);
    \item \texttt{heap} — динамическая память (\texttt{malloc}, \texttt{brk}), растёт вверх;
    \item \texttt{mmap} — отображаемые файлы и разделяемые библиотеки;
    \item \texttt{stack} — локальные переменные, параметры, адреса возврата, растёт вниз.
\end{enumerate}

\item Системные вызовы управления памятью \label{sec:4}

\begin{enumerate}

\item \texttt{mmap()}, \texttt{munmap()}

\texttt{mmap(addr, length, prot, flags, fd, offset)} — создаёт отображение
файла или анонимной памяти в адресное пространство процесса. Страницы
подгружаются по требованию (\emph{demand paging}).

\textbf{Флаги:}
\begin{enumerate}
    \item \texttt{MAP\_ANONYMOUS} — анонимное отображение (без файла);
    \item \texttt{MAP\_PRIVATE} — частное, с семантикой COW;
    \item \texttt{MAP\_SHARED} — общее, изменения видны другим процессам.
\end{enumerate}

\texttt{munmap(addr, length)} — удаляет отображение; адрес должен быть
выровнен по странице.

\item \texttt{mlock()}, \texttt{munlock()}, \texttt{mlockall()}, \texttt{munlockall()}

\texttt{mlock(addr, len)} — фиксирует страницы в RAM, запрещая выгрузку в
swap. Применяется в real-time приложениях и для криптографических ключей.

\texttt{munlock()} — снимает фиксацию; \texttt{mlockall(MCL\_CURRENT)} —
для всего адресного пространства.

\item \texttt{mprotect()}

\texttt{mprotect(addr, len, prot)} — изменяет права доступа к страницам:
\texttt{PROT\_READ}, \texttt{PROT\_WRITE}, \texttt{PROT\_EXEC},
\texttt{PROT\_NONE}.

\texttt{PROT\_NONE} используется для \emph{guard pages} — обращение к ним
приводит к \texttt{SIGSEGV}.

\item \texttt{madvise()}, \texttt{posix\_madvise()}

\texttt{madvise(addr, len, advice)} — передаёт ядру советы по работе с
диапазоном:
\begin{enumerate}
    \item \texttt{MADV\_SEQUENTIAL} — ожидается последовательный доступ;
    \item \texttt{MADV\_RANDOM} — случайный доступ, readahead не нужен;
    \item \texttt{MADV\_WILLNEED} — диапазон скоро потребуется;
    \item \texttt{MADV\_DONTNEED} — диапазон пока не нужен;
    \item \texttt{MADV\_HUGEPAGE} / \texttt{MADV\_NOHUGEPAGE} — THP;
    \item \texttt{MADV\_MERGEABLE} — пометить для KSM.
\end{enumerate}

\texttt{posix\_madvise()} — POSIX-аналог с другим набором констант.
\end{enumerate}

\item \texttt{malloc()} и системные вызовы \label{sec:5}

\texttt{malloc}, \texttt{calloc}, \texttt{realloc}, \texttt{free} — функции
\emph{пользовательского аллокатора} (glibc), а не системные вызовы.

\begin{enumerate}
    \item небольшие блоки — из внутреннего пула через \texttt{brk()}/\texttt{sbrk()};
    \item крупные (\(> \texttt{MMAP\_THRESHOLD}\), обычно 128~КБ) — через \texttt{mmap()};
    \item \texttt{free()} не возвращает память ядру сразу — блок остаётся в пуле.
\end{enumerate}


\item Механизмы работы ядра с памятью \label{sec:6}

\begin{enumerate}

\item Copy-on-Write (COW)

После \texttt{fork()} родитель и потомок разделяют физические страницы
(\emph{read-only}). При записи возникает page fault, ядро копирует
страницу только пишущему процессу.

COW также применяется в \texttt{mmap}, снапшотах, контейнерах, VM, KSM.

\item Подкачка (swap)

Swap — область на диске для временного хранения неактивных
\emph{анонимных} страниц.

\begin{enumerate}
    \item \textbf{swap out} — RAM \(\to\) swap;
    \item \textbf{swap in} — swap \(\to\) RAM.
\end{enumerate}

Параметр \texttt{vm.swappiness} (\(\sim\)60 по умолчанию) влияет на
склонность к вытеснению анонимных страниц.

\item Кэш страниц (page cache)

Page cache — кэш содержимого файлов в RAM. \texttt{read()} и файловый
\texttt{mmap()} используют одни и те же страницы.

\begin{enumerate}
    \item \textbf{cold cache} — требуется реальный I/O;
    \item \textbf{warm cache} — данные уже в RAM.
\end{enumerate}

Writeback — асинхронная запись <<грязных>> страниц; \texttt{msync()} —
принудительная синхронизация.

\item Page fault

\begin{enumerate}
    \item \textbf{minor} — страница уже в RAM, нужно обновить PTE;
    \item \textbf{major} — требуется дисковый I/O (файл или swap);
    \item \textbf{SIGSEGV} — ядро не может законно разрешить обращение.
\end{enumerate}

\textbf{Нормальные причины:} demand paging, COW, подключение страницы,
уже находящейся в RAM.
\end{enumerate}


\item Консольные инструменты \label{sec:7}

\begin{enumerate}

\item \texttt{free}

\texttt{free -h} — обзор использования RAM и swap.

\item \texttt{vmstat}

\texttt{vmstat 1} — динамика памяти, подкачки (\texttt{si}/\texttt{so}),
CPU и I/O.

\item \texttt{top} / \texttt{htop}

Интерактивный мониторинг процессов и их памяти (RSS, VSZ, CPU).

\item \texttt{/proc/[PID]/maps} и \texttt{/proc/[PID]/smaps}

\texttt{maps} — карта VMA (адреса, права, offset, inode, файл).
\texttt{smaps} — то же плюс RSS, PSS, Private\_Dirty, Shared\_Dirty.

\item \texttt{/proc/meminfo}

Сводная информация: MemTotal, MemFree, Buffers, Cached, SwapTotal,
Committed\_AS, CommitLimit.

\item \texttt{pmap}

\texttt{pmap -X <PID>} — наглядное отображение карт памяти процесса.

\end{enumerate}

\item Сводная таблица

\begin{center}
\begin{tabular}{|l|l|}
\hline
\textbf{Инструмент / вызов} & \textbf{Назначение} \\
\hline
\texttt{getconf PAGE\_SIZE}    & размер страницы \\
\texttt{/proc/<PID>/maps}      & VMA процесса \\
\texttt{/proc/<PID>/smaps}     & RSS, PSS, Dirty по VMA \\
\texttt{/proc/<PID>/status}    & VmSize, VmRSS, VmSwap \\
\texttt{/proc/meminfo}         & общая статистика памяти \\
\texttt{/proc/vmstat}          & pgfault, pgmajfault, pswpin/out \\
\texttt{ps -o min\_flt,maj\_flt} & page faults процесса \\
\texttt{perf stat -e \ldots}   & dTLB, cache, faults \\
\hline
\end{tabular}
\end{center}
\end{enumerate}